\subsubsection{Relación objetivos específicos y resultados esperados}

% avoid nested structures in table
\newcommand\objetivoEspecificoAResulados{
  \begin{itemize}[leftmargin=*, topsep=0pt, itemsep=0pt]
    \item Lógica crítica de los módulos de procesos y campañas ejecutándose en entornos del lado del servidor (AWS Lambda y BFF en Nuxt), sin exposición en el cliente.
    \item Integraciones con APIs externas (Rama Judicial y SAMAI) encapsuladas en el BFF como \textit{proxy} controlado, con saneamiento y manejo centralizado de errores.
  \end{itemize}
}

\newcommand\objetivoEspecificoBResultados{
  \begin{itemize}[leftmargin=*, topsep=0pt, itemsep=0pt]
    \item \textit{Frontend} organizado en Nuxt Layers por \textit{bounded context} bajo DDD, con tema visual aislado en una capa independiente.
    \item \textit{Backend} en Go con separación estricta \textit{Handler--Service--Repository} y repositorios definidos por interfaz, que habilitan pruebas unitarias sin acceso a la base de datos real y preservan la opción de sustituir el motor de persistencia sin tocar la lógica de negocio.
  \end{itemize}
}

\newcommand\objetivoEspecificoCResultados{
  \begin{itemize}[leftmargin=*, topsep=0pt, itemsep=0pt]
    \item Telemetría estructurada en formato JSON emitida por interceptores HTTP en el \textit{frontend} (BFF/Nuxt) y en el \textit{backend} (Go), con identificadores de correlación propagados.
    \item Dashboards de CloudWatch, alarmas configuradas y canal de Telegram operativos para la recepción de notificaciones proactivas ante anomalías.
    \item Infraestructura del subsistema de observabilidad definida y versionada como código mediante CloudFormation.
  \end{itemize}
}

\newcommand\objetivoEspecificoDResultados{
  \begin{itemize}[leftmargin=1em, topsep=2pt, itemsep=2pt]
    \item \textit{Hooks} de \textit{pre-commit}, \textit{commit-msg} y \textit{post-merge} operativos a través de Lefthook, integrando ESLint, Prettier y Commitlint.
    \item Pipeline de integración continua en GitHub Actions que ejecuta validación diferencial sobre los archivos modificados en cada \textit{pull request}.
    \item Suite de pruebas \textit{end-to-end} con Playwright y suite de pruebas unitarias de las reglas de seguridad de Firestore con Vitest, ambas ejecutándose sobre la \textit{Firebase Emulator Suite}.
  \end{itemize}
}

\newcommand\objetivosEspecificosYEsperadosCaption{Objetivos específicos y resultados esperados del proyecto. \hspace{1em}}

\begingroup
\renewcommand{\arraystretch}{1.2}
\begin{longtable}{|p{8cm}|p{8cm}|}
  \hline
  % Table Header
  \grayTableHeaderCell{8cm}{Objetivo específico} &
  \grayTableHeaderCell{8cm}{Resultados esperados} \\
  \hline

  % Table Body
  \tableCell{\objetivoEspecificoA}   &
  \tableCell{\objetivoEspecificoAResulados} \\
  \hline

  \tableCell{\objetivoEspecificoB}   &
  \tableCell{\objetivoEspecificoBResultados} \\
  \hline
  \endfirsthead

  \tableCell{\objetivoEspecificoB}   &
  \tableCell{\objetivoEspecificoBResultados} \\
  \hline
  \endhead

  \continuacionTablaFooter{2}
  \endfoot
  \endlastfoot

  \tableCell{\objetivoEspecificoC}   &
  \tableCell{\objetivoEspecificoCResultados} \\
  \hline

  \tableCell{\objetivoEspecificoD}   &
  \tableCell{\objetivoEspecificoDResultados} \\
  \hline
  \caption[\objetivosEspecificosYEsperadosCaption]{\objetivosEspecificosYEsperadosCaption Fuente: Elaboración propia.}
  \label{tab:objetivos_resultados}
\end{longtable}
\endgroup
